\documentclass[a4paper,landscape,8pt]{extarticle}
\usepackage[landscape,margin=0.7cm,top=0.5cm,bottom=0.7cm]{geometry}
\usepackage[utf8]{inputenc}
\usepackage[T1]{fontenc}
\usepackage[ngerman]{babel}
\usepackage{helvet}
\renewcommand{\familydefault}{\sfdefault}
\usepackage{xcolor}
\usepackage{tikz}
\usetikzlibrary{positioning}
\usepackage{enumitem}
\usepackage{multicol}
\usepackage{array}
\usepackage{fancyhdr}
\usepackage{amssymb}

% Farben
\definecolor{noteyellow}{HTML}{FFCC00}
\definecolor{dictblue}{HTML}{1565C0}
\definecolor{sharegreen}{HTML}{2E7D32}
\definecolor{lightgray}{HTML}{F5F5F5}
\definecolor{bordergray}{HTML}{BBBBBB}

\pagestyle{fancy}
\fancyhf{}
\renewcommand{\headrulewidth}{0pt}
\fancyfoot[C]{\footnotesize Seite \thepage{} von 3}

\setlength{\parindent}{0pt}
\setlength{\parskip}{2pt}

% Kompakte Listen
\setlist{nosep, leftmargin=1em, itemsep=0pt, topsep=1pt}

\begin{document}

% ============================================================================
% SEITE 1: INFORMATIONSBLATT
% ============================================================================

\noindent{\Large\bfseries Notizen: Dein digitaler Notizblock} \hfill
{\footnotesize\textbf{Lernziele:} Notizen erstellen | Diktieren | Ordnung halten}

\vspace{1mm}
\hrule
\vspace{1mm}

% Drei-Spalten-Layout
\begin{minipage}[t]{0.32\linewidth}
\begin{center}
\colorbox{noteyellow!30}{%
\begin{minipage}{0.95\linewidth}
\vspace{2mm}
\begin{center}
{\Large\bfseries\textcolor{noteyellow!80!black}{1. Die Notizen-App}}\\[1pt]
{\footnotesize\itshape Der digitale Zettel}
\end{center}
\vspace{2mm}
\end{minipage}%
}
\end{center}

\vspace{1mm}

\textbf{Was ist das?}\\
Ein digitaler Notizblock, der \textbf{nie voll wird} und den du immer dabei hast.

\textbf{So findest du sie:}
\begin{itemize}
\item Gelbes Symbol mit Stift
\item Auf dem Home-Bildschirm
\end{itemize}

\textbf{So erstellst du eine Notiz:}
\begin{enumerate}
\item App öffnen
\item Unten rechts: Quadrat mit Stift
\item Sofort losschreiben
\item Erste Zeile = Titel
\end{enumerate}

\textbf{Analogie:} Wie ein Notizblock -- nur besser, weil durchsuchbar!

\textbf{Praktische Ideen:}
\begin{itemize}
\item Einkaufsliste
\item Wichtige Telefonnummern
\item Medikamenten-Liste
\item Geburtstage
\end{itemize}

\end{minipage}%
\hfill%
\begin{minipage}[t]{0.32\linewidth}
\begin{center}
\colorbox{dictblue!10}{%
\begin{minipage}{0.95\linewidth}
\vspace{2mm}
\begin{center}
{\Large\bfseries\textcolor{dictblue}{2. Diktieren statt Tippen}}\\[1pt]
{\footnotesize\itshape Sprechen ist einfacher}
\end{center}
\vspace{2mm}
\end{minipage}%
}
\end{center}

\vspace{1mm}

\textbf{So geht's:}
\begin{enumerate}
\item Notiz öffnen
\item \textbf{Mikrofon-Symbol} auf der Tastatur antippen
\item Deutlich und langsam sprechen
\item Text erscheint automatisch
\item Nochmal antippen zum Beenden
\end{enumerate}

\textbf{Satzzeichen mitsprechen:}
\begin{itemize}
\item ``\textbf{Punkt}'' = .
\item ``\textbf{Komma}'' = ,
\item ``\textbf{Fragezeichen}'' = ?
\item ``\textbf{Neue Zeile}'' = Absatz
\end{itemize}

\textbf{Beispiel:}\\
``Einkaufsliste Doppelpunkt neue Zeile Milch Komma Brot Komma Äpfel Punkt''

\textbf{Ergebnis:}\\
Einkaufsliste:\\
Milch, Brot, Äpfel.

\vspace{2mm}
\fcolorbox{dictblue}{dictblue!5}{%
\begin{minipage}{0.92\linewidth}
\footnotesize
\textbf{Tipp:} Diktieren braucht Internet (WLAN oder mobile Daten)!
\end{minipage}%
}

\end{minipage}%
\hfill%
\begin{minipage}[t]{0.32\linewidth}
\begin{center}
\colorbox{sharegreen!10}{%
\begin{minipage}{0.95\linewidth}
\vspace{2mm}
\begin{center}
{\Large\bfseries\textcolor{sharegreen}{3. Ordnen \& Teilen}}\\[1pt]
{\footnotesize\itshape Übersicht behalten}
\end{center}
\vspace{2mm}
\end{minipage}%
}
\end{center}

\vspace{1mm}

\textbf{Notiz wiederfinden:}
\begin{itemize}
\item Oben: Suchfeld antippen
\item Stichwort eingeben
\item Alle passenden Notizen erscheinen
\end{itemize}

\textbf{Ordner erstellen:}
\begin{enumerate}
\item In der Ordner-Übersicht (links)
\item Unten links: ``Neuer Ordner''
\item Namen eingeben (z.B. ``Gesundheit'')
\end{enumerate}

\textbf{Notiz verschieben:}
\begin{enumerate}
\item Notiz antippen und halten
\item ``Bewegen'' wählen
\item Ziel-Ordner antippen
\end{enumerate}

\textbf{Notiz teilen:}
\begin{enumerate}
\item Notiz öffnen
\item Teilen-Symbol (Quadrat mit Pfeil)
\item Per Mail, WhatsApp senden
\end{enumerate}

\textbf{Ordner-Ideen:}
\begin{itemize}
\item Einkauf
\item Gesundheit
\item Familie
\item Reisen
\end{itemize}

\end{minipage}

\vspace{1mm}

% Zusammenfassung
\begin{center}
\fcolorbox{bordergray}{lightgray}{%
\begin{minipage}{0.98\linewidth}
\centering
\vspace{1.5mm}
\textbf{Zusammenfassung:}
\textcolor{noteyellow!80!black}{\textbf{Neue Notiz}} = Quadrat mit Stift ~$|$~
\textcolor{dictblue}{\textbf{Diktieren}} = Mikrofon-Symbol ~$|$~
\textcolor{sharegreen}{\textbf{Suchen}} = oben im Suchfeld ~$|$~
\textbf{Teilen} = Quadrat mit Pfeil
\vspace{1.5mm}
\end{minipage}%
}
\end{center}

\newpage

% ============================================================================
% SEITE 2: ÜBUNGEN
% ============================================================================

\begin{center}
{\LARGE\bfseries Übungen: Notizen meistern}\\[3pt]
{\normalsize Schau dir Seite 1 an und beantworte die Fragen}
\end{center}

\vspace{3mm}
\hrule
\vspace{4mm}

\begin{multicols}{2}

\textbf{\large Teil A: Grundlagen}

Richtig oder Falsch? Kreuze an.

\vspace{2mm}

\renewcommand{\arraystretch}{1.6}
\begin{tabular}{@{}p{8cm}cc@{}}
& R & F \\
1. Die Notizen-App hat ein gelbes Symbol. & $\square$ & $\square$ \\
2. Die erste Zeile wird automatisch zum Titel. & $\square$ & $\square$ \\
3. Diktieren funktioniert ohne Internet. & $\square$ & $\square$ \\
4. Notizen können in Ordnern sortiert werden. & $\square$ & $\square$ \\
5. Eine Notiz kann nur per E-Mail geteilt werden. & $\square$ & $\square$ \\
\end{tabular}
\renewcommand{\arraystretch}{1}

\vspace{6mm}

\textbf{\large Teil B: Diktieren}

Was musst du sagen, um diese Satzzeichen zu bekommen?

\vspace{2mm}

\renewcommand{\arraystretch}{1.8}
\begin{tabular}{@{}p{4cm}p{5cm}@{}}
1. Einen Punkt: & \dotfill \\
2. Ein Komma: & \dotfill \\
3. Ein Fragezeichen: & \dotfill \\
4. Einen Absatz: & \dotfill \\
\end{tabular}
\renewcommand{\arraystretch}{1}

\vspace{6mm}

\textbf{\large Teil C: Reihenfolge}

Nummeriere die Schritte zum Erstellen einer Notiz (1--4):

\vspace{2mm}

\renewcommand{\arraystretch}{1.8}
\begin{tabular}{@{}cp{7cm}@{}}
$\square$ & Text eingeben (tippen oder diktieren) \\
$\square$ & Quadrat mit Stift antippen \\
$\square$ & Notizen-App öffnen \\
$\square$ & Fertig -- Notiz speichert automatisch \\
\end{tabular}
\renewcommand{\arraystretch}{1}

\columnbreak

\textbf{\large Teil D: Praktisch üben}

Erstelle diese Notizen auf deinem Gerät:

\vspace{3mm}

\textbf{Aufgabe 1: Einkaufsliste}\\
Erstelle eine Notiz mit dem Titel ``Einkaufsliste'' und mindestens 5 Artikeln.

\vspace{2mm}
$\square$ Erledigt

\vspace{5mm}

\textbf{Aufgabe 2: Diktieren testen}\\
Erstelle eine neue Notiz und diktiere: ``Das ist meine erste diktierte Notiz Punkt''

\vspace{2mm}
$\square$ Erledigt

\vspace{5mm}

\textbf{Aufgabe 3: Notiz suchen}\\
Nutze die Suche, um deine Einkaufsliste wiederzufinden.

\vspace{2mm}
$\square$ Erledigt

\vspace{5mm}

\textbf{Aufgabe 4: Ordner erstellen}\\
Erstelle einen Ordner namens ``Privat'' oder ``Einkauf''.

\vspace{2mm}
$\square$ Erledigt

\vspace{5mm}

\textbf{Bonus: Notiz teilen}\\
Teile deine Einkaufsliste per WhatsApp oder E-Mail.

\vspace{2mm}
$\square$ Erledigt

\end{multicols}

\newpage

% ============================================================================
% SEITE 3: CHECKLISTE
% ============================================================================

\begin{center}
{\LARGE\bfseries Checkliste: Das habe ich verstanden}\\[3pt]
{\normalsize Geh die Liste durch und hake ab, was du sicher weißt}
\end{center}

\vspace{3mm}
\hrule
\vspace{6mm}

% Drei Boxen nebeneinander
\begin{minipage}[t]{0.31\linewidth}
\begin{center}
\fcolorbox{noteyellow!80!black}{noteyellow!20}{%
\begin{minipage}{0.92\linewidth}
\vspace{3mm}
\begin{center}
{\large\bfseries\textcolor{noteyellow!80!black}{Notizen erstellen}}
\end{center}
\vspace{2mm}
\begin{itemize}[label=$\square$, itemsep=5pt]
\item Ich finde die \textbf{Notizen-App}
\item Ich kann eine \textbf{neue Notiz} erstellen
\item Ich weiß, dass die erste Zeile der \textbf{Titel} wird
\item Notizen werden \textbf{automatisch} gespeichert
\end{itemize}
\vspace{3mm}
\end{minipage}%
}
\end{center}
\end{minipage}%
\hfill%
\begin{minipage}[t]{0.31\linewidth}
\begin{center}
\fcolorbox{dictblue}{dictblue!8}{%
\begin{minipage}{0.92\linewidth}
\vspace{3mm}
\begin{center}
{\large\bfseries\textcolor{dictblue}{Diktieren}}
\end{center}
\vspace{2mm}
\begin{itemize}[label=$\square$, itemsep=5pt]
\item Ich finde das \textbf{Mikrofon-Symbol}
\item Ich spreche \textbf{Satzzeichen} mit
\item Ich weiß: Diktieren braucht \textbf{Internet}
\item Ich kann eine Notiz \textbf{diktieren}
\end{itemize}
\vspace{3mm}
\end{minipage}%
}
\end{center}
\end{minipage}%
\hfill%
\begin{minipage}[t]{0.31\linewidth}
\begin{center}
\fcolorbox{sharegreen}{sharegreen!8}{%
\begin{minipage}{0.92\linewidth}
\vspace{3mm}
\begin{center}
{\large\bfseries\textcolor{sharegreen}{Ordnen \& Teilen}}
\end{center}
\vspace{2mm}
\begin{itemize}[label=$\square$, itemsep=5pt]
\item Ich kann Notizen \textbf{suchen}
\item Ich kann \textbf{Ordner} erstellen
\item Ich kann Notizen \textbf{verschieben}
\item Ich kann eine Notiz \textbf{teilen}
\end{itemize}
\vspace{3mm}
\end{minipage}%
}
\end{center}
\end{minipage}

\vspace{8mm}

% Gesamtverständnis
\begin{center}
\fcolorbox{black}{lightgray}{%
\begin{minipage}{0.95\linewidth}
\vspace{4mm}
\begin{center}
{\large\bfseries Gesamtverständnis}
\end{center}
\vspace{3mm}
\begin{multicols}{2}
\begin{itemize}[label=$\square$, itemsep=6pt]
\item Ich kann \textbf{Notizen erstellen} statt Zettel zu schreiben
\item Ich kann \textbf{diktieren}, wenn Tippen schwerfällt
\item Ich finde meine Notizen durch \textbf{Suchen} wieder
\item Ich halte \textbf{Ordnung} mit Ordnern
\item Ich kann Notizen \textbf{teilen} (Einkaufsliste senden)
\end{itemize}
\end{multicols}
\vspace{3mm}
\end{minipage}%
}
\end{center}

\vspace{8mm}

% Notfall-Notizen
\begin{center}
\fcolorbox{bordergray}{white}{%
\begin{minipage}{0.95\linewidth}
\vspace{4mm}
\begin{center}
{\large\bfseries Meine Notiz-Ideen}\\[3pt]
{\footnotesize (Was möchtest du in Notizen festhalten?)}
\end{center}
\vspace{4mm}
\renewcommand{\arraystretch}{2}
\begin{tabular}{@{}p{0.45\linewidth}p{0.45\linewidth}@{}}
\textbf{Praktische Notizen:} & \textbf{Meine Ordner:} \\
$\square$ Einkaufsliste & 1. \dotfill \\[2mm]
$\square$ Medikamente & 2. \dotfill \\[2mm]
$\square$ Wichtige Telefonnummern & 3. \dotfill \\[2mm]
$\square$ Geburtstage & \\
$\square$ \dotfill & \\
\end{tabular}
\renewcommand{\arraystretch}{1}
\vspace{4mm}
\end{minipage}%
}
\end{center}

\vfill

\begin{center}
\footnotesize\textit{Stand: \today{} -- Notizen ersetzen Zettelwirtschaft -- immer dabei, nie verloren!}
\end{center}

\end{document}
